\chapter{GÓI PHÁT HÀNH VÀ HƯỚNG DẪN SỬ DỤNG}

\section{Deliverable Package}

\begin{longtable}{|c|p{4cm}|p{8cm}|}
	\hline
	\textbf{No.} & \textbf{Deliverable Item} & \textbf{Description} \\ 
	\hline
	1 & Schedule/Task Tracking &
	Sử dụng Jira để theo dõi tiến độ thực hiện các công việc của dự án và quản lý các task của từng thành viên. \\
	\hline
	
	2 & Project Backlog &
	Sử dụng Jira để quản lý Product Backlog, phân chia các chức năng thành các task và theo dõi trạng thái thực hiện của từng task. \\
	\hline
	
	3 & Source Codes &
	Mã nguồn của hệ thống được quản lý và lưu trữ trên GitHub, bao gồm:
	\begin{itemize}
		\item Back End: Flask Clean Architecture
		\item Front End: Web application sử dụng HTML, CSS và JavaScript
	\end{itemize}
	\\
	\hline
	
	4 & Database Script(s) &
	Các script SQL được sử dụng để tạo và quản lý cơ sở dữ liệu PostgreSQL của hệ thống. Cơ sở dữ liệu được triển khai trên Supabase. \\
	\hline
	
	5 & Final Report Document &
	Tài liệu báo cáo cuối kỳ tổng hợp các nội dung của dự án, bao gồm phân tích yêu cầu, thiết kế hệ thống, cơ sở dữ liệu, kiểm thử và hướng dẫn sử dụng. \\
	\hline
	
	6 & Test Cases Document &
	Tài liệu mô tả phạm vi, chiến lược và kế hoạch kiểm thử cho các chức năng của hệ thống. \\
	\hline
	
	7 & Slide &
	Tài liệu trình chiếu giới thiệu dự án, bao gồm thành viên nhóm, actors, bối cảnh dự án, công nghệ sử dụng, yêu cầu chức năng, yêu cầu phi chức năng và các luồng chức năng chính của hệ thống. \\
	\hline
	
\end{longtable}

\section{Hướng dẫn cài đặt}

\subsection{Yêu cầu hệ thống}

\subsubsection{Web Application}

\begin{longtable}{|p{4cm}|p{5cm}|p{5cm}|}
	\hline
	\textbf{Component} & \textbf{Minimum} & \textbf{Recommended} \\
	\hline
	CPU & Intel Core i5 & Intel Core i7 \\
	\hline
	Memory & 8GB RAM & 16GB RAM trở lên \\
	\hline
	Storage Space & 10GB trống & 20GB trống \\
	\hline
	Network & 10 Mbps & 50 Mbps trở lên \\
	\hline
	Web Browser & Chrome phiên bản mới & Chrome phiên bản mới nhất \\
	\hline
\end{longtable}

\subsection{Hướng dẫn cài đặt}

\subsubsection{Backend}

\textbf{Step 1:} Cài đặt Visual Studio Code.

\textbf{Step 2:} Cài đặt Git.

\textbf{Step 3:} Clone source code Backend từ GitHub.

\textbf{Step 4:} Mở thư mục Backend bằng Visual Studio Code.

\textbf{Step 5:} Tạo và kích hoạt môi trường ảo Python.

\textbf{Step 6:} Cài đặt các thư viện cần thiết.

\textbf{Step 7:} Cấu hình thông tin kết nối cơ sở dữ liệu PostgreSQL.

\textbf{Step 8:} Chạy Backend.

\subsubsection{Frontend}

\textbf{Step 1:} Mở thư mục Frontend bằng Visual Studio Code.

\textbf{Step 2:} Kiểm tra cấu hình kết nối với Backend.

\textbf{Step 3:} Mở Terminal tại thư mục Frontend.

\textbf{Step 4:} Chạy Frontend bằng trình duyệt web.

\textbf{Step 5:} Truy cập địa chỉ của ứng dụng trên trình duyệt.
\section{Hướng dẫn sử dụng}

\subsection{Tổng quan}

% Liệt kê chức năng theo Photographer / Service Provider / Administrator

\subsection{Hướng dẫn chức năng}

\subsubsection{Đăng ký tài khoản}